\bigskip

% ─────────────────────────────────────────────
% Pages 15–21: Off-Policy Actor-Critic Methods
% ─────────────────────────────────────────────

\textbf{Page 15 --- The Big Question: Can We Reuse Old Data?}

\medskip
So far, the actor-critic algorithm only learns from data it just collected.
Once you are done with a batch of experience, you throw it away. That is wasteful.
Page 15 asks: What if we saved every single experience the agent ever had,
and learned from ALL of it --- not just the latest batch?

\medskip
\textbf{Two key ideas introduced:}
\begin{enumerate}
    \item \textbf{Replay Buffer} --- a giant diary that stores every
          (state, action, reward, next-state) experience the agent has ever seen.
    \item \textbf{Fix the equations} --- the original equations assume you only
          use fresh data. If you use old data, you need to adjust the math.
\end{enumerate}

\bigskip

\textbf{Page 16 --- First Attempt: Plug the Replay Buffer In}

\medskip
They take the standard algorithm and simply say: instead of doing steps on one
fresh batch, do them on a small random sample (``minibatch'') pulled from the
replay buffer --- which contains all past experience.

\medskip
\textbf{The 5-step algorithm (first attempt):}
\begin{enumerate}
    \item Collect fresh experience and add it to the replay buffer.
    \item Fit the value function $\hat{V}^{\pi}_{\phi}$ to reward sums from the buffer.
    \item Calculate how good each (state, action) was using the advantage:
          \[
              \hat{A}^{\pi}(s_i, a_i) = r(s_i, a_i) + \hat{V}^{\pi}_{\phi}(s'_i) - \hat{V}^{\pi}_{\phi}(s_i)
          \]
    \item Estimate the policy gradient.
    \item Update the policy using that gradient.
\end{enumerate}

Steps 2--4 are done on a minibatch randomly sampled from all previous data ---
that is the key change.

\bigskip

\textbf{Page 17 --- Two Things Are Broken}

\medskip
Once they lay out the algorithm, they realize there is a serious problem.
When $V^{\pi}$ is fit on ALL data from the replay buffer (which comes from many
different past versions of the policy), the value function represents some kind
of average over old policies --- not the current policy. This is a problem
because the value function is supposed to tell you how good the \textit{current}
policy is.

\medskip
\textbf{Two things broken:}
\begin{itemize}
    \item \textbf{The target is wrong.} The value estimate $y = r + \gamma \hat{V}(s')$
          uses a $\hat{V}$ trained on a mix of old policies, so it does not correctly
          represent the current policy's future rewards.
    \item \textbf{The likelihood is wrong.} In the gradient step, you are using the
          log probability of actions that old versions of the policy took --- not
          actions the current policy would take.
\end{itemize}

\bigskip

\textbf{Page 18 --- The Fix: Use $Q(s, a)$ Instead of $V(s)$}

\medskip
\textbf{Core insight:} Instead of estimating how good a state is ($V$), estimate
how good a state + action pair is ($Q$). This matters because $Q(s, a)$ can be
computed from any experience, regardless of which policy collected it.

\medskip
\textbf{The equation:}
\[
    Q^{\pi_\theta}(s, a)
    = r(s, a)
    + \gamma \,
      \mathbb{E}_{\substack{s' \sim p(\cdot|s,a) \\ \bar{a}' \sim \pi_\theta(\cdot|s')}}
      \!\left[ Q^{\pi_\theta}(s', \bar{a}') \right]
\]

\medskip
\textbf{Symbol table:}

\begin{tabular}{p{5.5cm} p{8.5cm}}
\hline
\textbf{Symbol} & \textbf{What it means} \\
\hline
$Q^{\pi_\theta}(s, a)$ & Q-value: expected total future reward starting in state $s$, taking action $a$, then following policy $\pi_\theta$ \\[4pt]
$r(s, a)$ & Immediate reward for taking action $a$ in state $s$ \\[4pt]
$\gamma$ & Discount factor (0 to 1); sooner rewards are worth more than later ones \\[4pt]
$\mathbb{E}$ & Expected value --- average over many possible outcomes \\[4pt]
$s' \sim p(\cdot \mid s, a)$ & Next state sampled from the environment's transition function \\[4pt]
$\bar{a}' \sim \pi_\theta(\cdot \mid s')$ & Next action sampled from the \textit{current} policy at next state $s'$ \\[4pt]
$Q^{\pi_\theta}(s', \bar{a}')$ & Q-value of the next state-action pair \\
\hline
\end{tabular}

\medskip
The quality of taking action $a$ in state $s$ equals the immediate reward right
now, plus the discounted quality of whatever state and action comes next.
The ``next action'' $\bar{a}'$ must be sampled from the \textit{current} policy ---
not retrieved from the replay buffer --- because we want to know how good the
current policy is, not some past version.
The approach: sample $(s, a, s')$ from the replay buffer (old data is fine),
but sample $\bar{a}'$ fresh from the current policy.
Then the target becomes $y = r(s, a) + \gamma Q(s', \bar{a}')$, and you
train the Q-network toward this target.

\bigskip

\textbf{Page 19 --- Fixing the Value Function: The Updated Algorithm}

\medskip
This page rewrites the algorithm with the fix from page 18 applied.

\medskip
\textbf{The key fix in step 3:}
\begin{itemize}
    \item \textbf{Old (broken):} Update $\hat{V}^{\pi}_\phi$ using targets
          $y_i = r_i + \gamma \hat{V}^{\pi}_\phi(s'_i)$
    \item \textbf{New (fixed):} Update $\hat{Q}^{\pi}_\phi$ using targets
          $y_i = r(s_i, a_i) + \gamma \hat{Q}^{\pi}_\phi(s'_i, a'_i)$,
          where $a'_i \sim \pi_\theta(\cdot \mid s'_i)$ --- sampled from the
          \textit{current} policy, not from the replay buffer.
\end{itemize}

\medskip
Using a fresh action fixes the broken target problem --- the Q-value target now
reflects the current policy. The states $s_i$ can still come from the replay buffer;
only the next actions need to be fresh.

\bigskip

\textbf{Page 20 --- Two More Refinements}

\medskip
\textbf{Refinement 1: Drop the baseline, just use $Q$ directly.}

Instead of computing the advantage $\hat{A}(s, a) = Q(s, a) - V(s)$, just use
$Q(s, a)$ directly. This has higher variance (noisier estimates), but since you
now have a huge replay buffer full of data, the extra noise is acceptable ---
more data cancels it out.

\medskip
\textbf{Refinement 2: Sample fresh actions for the gradient step too.}

\[
    \nabla_\theta J(\theta)
    \approx \frac{1}{N} \sum_i
    \nabla_\theta \log \pi_\theta(\bar{a}_i \mid s_i) \,
    \hat{Q}^{\pi}(s_i, \bar{a}_i)
\]

\medskip
\textbf{Symbol table:}

\begin{tabular}{p{5.5cm} p{8.5cm}}
\hline
\textbf{Symbol} & \textbf{What it means} \\
\hline
$\nabla_\theta J(\theta)$ & Gradient: the direction to nudge policy parameters $\theta$ to improve total reward \\[4pt]
$\frac{1}{N} \sum_i$ & Average over $N$ sampled examples \\[4pt]
$\nabla_\theta \log \pi_\theta(\bar{a}_i \mid s_i)$ & How much the log-probability of action $\bar{a}_i$ changes when you tweak $\theta$ \\[4pt]
$\bar{a}_i$ & Fresh action sampled from current policy $\pi_\theta$ at state $s_i$ --- NOT the old stored action \\[4pt]
$\hat{Q}^{\pi}(s_i, \bar{a}_i)$ & Estimated quality of taking fresh action $\bar{a}_i$ in state $s_i$ \\[4pt]
$\pi_\theta(\bar{a}_i \mid s_i)$ & Probability the current policy assigns to action $\bar{a}_i$ given state $s_i$ \\
\hline
\end{tabular}

\medskip
To improve the policy, look at a bunch of state-action pairs. For each state $s_i$
from the buffer, sample a fresh action $\bar{a}_i$ from the current policy (not the
old stored action). Check how good that fresh action is according to $\hat{Q}$.
Then nudge the policy parameters in the direction that makes good actions more likely.
Using a fresh action fixes the ``wrong likelihood'' problem from page 17 --- you are
now computing log probabilities under the current policy, not some old one.
The states $s_i$ can still come from the replay buffer; only the actions need to be fresh.

\bigskip

\textbf{Page 21 --- The Final Clean Algorithm}

\medskip
This is the polished, fixed version of the off-policy actor-critic with a replay buffer.

\medskip
\textbf{Final 5-step algorithm:}
\begin{enumerate}
    \item Collect experience $\{s_i, a_i\}$ from the current policy and add it
          to the replay buffer.
    \item Sample a batch $\{s_i, a_i, r_i, s'_i\}$ from the replay buffer.
    \item Update $\hat{Q}^{\pi}_\phi$ toward targets
          $y_i = r(s_i, a_i) + \gamma \hat{Q}^{\pi}_\phi(s'_i, a'_i)$,
          where $a'_i \sim \pi_\theta(\cdot \mid s'_i)$ --- sampled fresh from
          the current policy.
    \item Compute the gradient:
          \[
              \nabla_\theta J(\theta) \approx \frac{1}{N} \sum_i
              \nabla_\theta \log \pi_\theta(\bar{a}_i \mid s_i) \,
              \hat{Q}^{\pi}(s_i, \bar{a}_i)
          \]
          where $\bar{a}_i \sim \pi_\theta(\cdot \mid s_i)$ --- freshly sampled
          from the current policy.
    \item Update $\theta \leftarrow \theta + \alpha \nabla_\theta J(\theta)$.
\end{enumerate}

\medskip
The replay buffer lets the agent learn from every experience it has ever had,
making training much more data-efficient. The two key fixes --- sampling fresh
actions for the Q-target and for the gradient --- ensure that even though the
\textit{data} is old, the \textit{evaluation} always reflects what the current
policy would do.

\bigskip
\bigskip

% ─────────────────────────────────────────────
% Follow-up Questions
% ─────────────────────────────────────────────

\textbf{What is meant by random sample (minibatch)?}

\medskip
The replay buffer is like a giant jar full of thousands of past experiences
(state, action, reward, next-state). Instead of learning from every single one
of those experiences at once (too slow and memory-heavy), you randomly grab a
small handful --- say, 64 or 256 --- and learn from just those.

That small random handful is the \textbf{minibatch}.

``Random'' is important here: you do not pick the most recent ones or the best
ones --- you grab them randomly so the batch is not biased toward any one policy
or time period. This also breaks correlations between nearby experiences, which
helps the learning stay stable.

\bigskip

\textbf{Q(s,a) can be computed from any experience. How and why?}

\medskip
\textbf{First, recall what each one asks:}
\begin{itemize}
    \item $V(s)$: ``How good is this state, if I follow my current policy from
          here?'' --- depends entirely on the policy.
    \item $Q(s, a)$: ``How good is it to take action $a$ in state $s$, and then
          follow my current policy from there?'' --- the action is already pinned down.
\end{itemize}

\medskip
\textbf{Why $Q$ can use old data:}

The target for $Q$ is:
\[
    y = r(s, a) + \gamma \hat{Q}^{\pi}(s', \bar{a}')
\]

To compute this, you need three things:

\begin{tabular}{p{3cm} p{5cm} p{5cm}}
\hline
\textbf{Thing needed} & \textbf{Where it comes from} & \textbf{Policy-dependent?} \\
\hline
$r(s, a)$ & Replay buffer (old data) & No --- the environment gave this reward, not the policy \\[4pt]
$s'$ & Replay buffer (old data) & No --- the environment transitioned to $s'$, not the policy \\[4pt]
$\bar{a}'$ & Sampled fresh from current policy & Yes --- and you do this fresh \\
\hline
\end{tabular}

\medskip
The reward $r(s, a)$ is a fact about the environment --- it does not matter which
policy decided to take action $a$. If you step off a cliff in a game, you lose a
life whether your old policy did it or your new one. The environment always gives
the same reward for the same $(s, a)$. The next state $s'$ is also determined by
the environment, not the policy. So the $(s, a, r, s')$ tuple from old experience
is completely valid. The only policy-sensitive piece --- what action comes next ---
is handled by sampling $\bar{a}'$ fresh from the current policy at $s'$.

\medskip
\textbf{Why $V$ cannot use old data:}

$V(s)$ has no action pinned. It averages over whatever actions the policy would
take. So if your replay buffer has data from 10 different old policies, $V(s)$
ends up representing some confused blend of all of them --- not your current policy.
There is no way to ``correct'' for this without knowing exactly which old policy
produced each experience.

\medskip
\textbf{In short:} $Q$ separates the policy-independent part ($r$, $s'$) from the
policy-dependent part ($\bar{a}'$). Old data handles the first part; the current
policy handles the second. $V$ has no such clean separation.